\hypertarget{cont2:paper-ix-correction-torsors-and-exact-polynomial-lift-counts}{%
\subsection{Paper IX: correction torsors and exact polynomial lift
counts}\label{cont2:paper-ix-correction-torsors-and-exact-polynomial-lift-counts}}

\textbf{Author: Komorebi P. --- 8 September 2026.}

This block formalizes the correction-digit torsor and prime-field counts
from the previous checklist. The counts apply to distinct next-level
residue classes satisfying the actual polynomial equations. They do not
count integer representatives, of which there can be infinitely many.

\hypertarget{cont2:the-theorem}{%
\subsubsection{The theorem}\label{cont2:the-theorem}}

Let \(p\) be prime, \(L\ge1\), and \[
F\in\mathbb Z[X_1,\ldots,X_d]^e,
\qquad x\in\mathbb Z^d,
\qquad F(x)\equiv0\pmod{p^L}.
\] Write \[
r_x=F(x)/p^L,\qquad b_x=r_x\bmod p,
\qquad A_x=DF(x)\bmod p:\mathbb F_p^d\longrightarrow\mathbb F_p^e,
\qquad \rho_x=\operatorname{rank}_{\mathbb F_p}A_x.
\] The actual lift set is \[
\mathcal L_{x,L}=
\left\{y\in(\mathbb Z/p^{L+1}\mathbb Z)^d:
 y\equiv x\pmod{p^L},\quad F(y)=0\right\},
\] where polynomial evaluation in the displayed set is modulo
\(p^{L+1}\). Define \[
\mathcal O_x=\mathbb F_p^e/\operatorname{im}A_x,
\qquad o_x=[b_x].
\] Lean proves \[
\boxed{
\#\mathcal L_{x,L}=
\begin{cases}
p^{d-\rho_x},&o_x=0,\\
0,&o_x\ne0,
\end{cases}
\qquad
\#\mathcal O_x=p^{e-\rho_x}.
}
\] These formulas permit singular Jacobians and zero variables or zero
equations. Nonsingularity is not assumed. In particular, rank alone
gives the size of a nonempty lift fiber; the residual obstruction
determines whether it is empty.

\hypertarget{cont2:why-these-are-actual-residue-class-counts}{%
\subsubsection{Why these are actual residue-class
counts}\label{cont2:why-these-are-actual-residue-class-counts}}

For every positive integer \(q\), including composite moduli and modulus
one, and every \(L\ge0\), the checked map \[
\Phi_{x,L}:(\mathbb Z/q\mathbb Z)^d\longrightarrow
\{y\bmod q^{L+1}:y\equiv x\pmod{q^L}\},
\qquad t\longmapsto x+q^L\widetilde t
\] is a bijection. Here \(\widetilde t\) is the canonical integer
representative of the digit vector. The entire residue fiber therefore
has \(q^d\) elements. Injectivity is proved by cancelling \(q^L\) in the
corresponding integer divisibility statement; surjectivity comes from
the lower-residue congruence.

The formal type \mcode{ActualLifts} is a subtype of these genuine
\texttt{ZMod\ (q\^{}(L+1))} residue vectors. Its polynomial predicate is
evaluated on canonical integer representatives and tested for
divisibility by \(q^{L+1}\). The checked polynomial congruence theorem
makes that predicate independent of the integer representative. Thus
using canonical representatives does not discard or duplicate residue
classes.

At positive depth, the preceding actual multivariable Taylor theorem
gives \[
F(x+p^L\widetilde t)\equiv
p^L(r_x+DF(x)\widetilde t)\pmod{p^{L+1}}.
\] The new bridge proves that the evaluated integer Jacobian reduces to
the actual \(\mathbb F_p\)-linear map \(A_x\). Restricting
\(\Phi_{x,L}\) therefore gives an explicit equivalence \[
\boxed{\{t\in\mathbb F_p^d:b_x+A_xt=0\}\simeq\mathcal L_{x,L}.}
\] The proof uses the actual formal partial derivatives of the original
polynomial list.

\hypertarget{cont2:the-torsor}{%
\subsubsection{The torsor}\label{cont2:the-torsor}}

Let \(\mathcal T_x=\{t:b_x+A_xt=0\}\). Its natural kernel action is \[
k\cdot t=k+t,\qquad k\in\ker A_x.
\] For every pair \(u,t\in\mathcal T_x\), exactly one kernel vector
sends \(u\) to \(t\), namely \(t-u\). Thus every nonempty correction
fiber is a torsor for \(\ker A_x\). Choosing one valid digit \(u\) gives
\[
\mathcal T_x\simeq\ker A_x,
\qquad t\longmapsto t-u.
\] Lean constructs an explicit \mcode{AddTorsor} structure for the
digit fiber, proves the unique translation property, and supplies a
chosen-point equivalence from the actual residue lift set to the kernel.
A basepoint is required for this equivalence, but the translation action
itself does not depend on choosing one.

The torsor and residue bijection remain valid over every positive
modulus at positive polynomial depth. The displayed powers involving
vector-space rank are restricted to prime moduli.

\hypertarget{cont2:connecting-the-old-obstruction-to-the-new-count}{%
\subsubsection{Connecting the old obstruction to the new
count}\label{cont2:connecting-the-old-obstruction-to-the-new-count}}

The earlier obstruction used the integer-domain map \[
J_x=\rho_p\circ DF(x):\mathbb Z^d\to\mathbb F_p^e.
\] The new count uses \(A_x:\mathbb F_p^d\to\mathbb F_p^e\). Reduction
of integer vectors onto finite vectors is surjective, so the two maps
have the same image after restricting scalars. Lean constructs a
class-preserving integer-linear equivalence \[
\mathbb F_p^e/\operatorname{im}J_x
\simeq
\mathbb F_p^e/\operatorname{im}A_x.
\] Consequently the cardinality theorem counts the original obstruction
group, preserving its distinguished residual class. It is not merely a
formula for a separately introduced comparison object.

Rank-nullity gives \(\dim\ker A_x=d-\rho_x\), and the quotient dimension
is \(e-\rho_x\). Cardinalities of finite vector spaces complete both
formulas.
